% ===============================
% MAT221 -- Real Analysis
% Riemann Integral
% ===============================

\documentclass[12pt,a4paper]{article}

% ===============================
% Metadata
% ===============================
\newcommand*{\CourseCode}{MAT221}
\newcommand*{\CourseTitle}{Real Analysis}
\newcommand*{\Chapter}{Riemann Integrability}
\newcommand*{\Topics}{Integrating Functions with Discontinuities}
\newcommand*{\DocDate}{August 18, 2026}

% ===============================
% Header
% ===============================
\input{header}

% ===============================
% Document
% ===============================
\begin{document}

\FrontPage
\Large

\begin{heading}
    Continuous Functions are Riemann Integrable
\end{heading}

\vspace{10pt}

\begin{theorem}[Continuous Functions are Riemann Integrable]
    Let $f$ be a continuous function on a closed interval $[a,b]$. Then $f$ is Riemann integrable on $[a,b]$.
\end{theorem}

\clearpage\begin{heading}
    Monotone Functions are Riemann Integrable
\end{heading}

\vspace{10pt}

\begin{theorem}[Monotone Functions are Riemann Integrable]
    Let $f$ be a monotone function on a closed interval $[a,b]$. Then $f$ is Riemann integrable on $[a,b]$.
\end{theorem}

\clearpage


\begin{heading}
    Discontinuous Function May Not Be Riemann Integrable
\end{heading}

\vspace{10pt}

\begin{problem}[Dirichlet Function]
    Let $f$ be the Dirichlet function defined on the interval $[0,1]$ by
    \[
        f(x) = 
        \begin{cases}
            1, & \text{if } x \in \mathbb{Q},\\
            0, & \text{if } x \notin  \mathbb{Q}.
        \end{cases}
    \]
    Show that $f$ is not Riemann integrable on $[0,1]$.
    
\end{problem}

\clearpage

\begin{heading}
    Some Discontinuous Functions are Riemann Integrable
\end{heading}


\begin{theorem}[Discontinuity at one endpoint]
    Let $f$ be the function defined on the interval $[0,1]$ by
    \[
        f(x) = 
        \begin{cases}
            1, & \text{if } 0 < x < 1,\\
            0, & \text{if } x = 1.
        \end{cases}
    \]
    Show that $f$ is Riemann integrable on $[0,1]$ and compute its integral.
\end{theorem}

\clearpage

\begin{theorem}[Discontinuity at a point]
    Let $f$ be the function defined on the interval $[0,2]$ by
    \[
        f(x) = 
        \begin{cases}
            1, & \text{if } x \ne 1,\\
            0, & \text{if } x = 1.
        \end{cases}
    \]
    Show that $f$ is Riemann integrable on $[0,2]$ and compute its integral.
\end{theorem}

\clearpage

\begin{theorem}[Changing the value of a function at a point]
    Assume $f: [a,b] \to \mathbb{R}$ is Riemann integrable on $[a,b]$. Show that if one value of $f$ is changed at a single point $c \in [a,b]$, then the new function is also Riemann integrable on $[a,b]$ and the value of the integral does not change.
\end{theorem}

\clearpage

\begin{heading}
    Functions with Finite Number of Discontinuities
\end{heading}

\vspace{10pt}

\begin{theorem}
    Let $f$ be a bounded function on a closed interval $[a,b]$ that has only a finite number of discontinuities. Then $f$ is Riemann integrable on $[a,b]$.
    
\end{theorem}

\clearpage

\begin{theorem}[Changing value at finite number of points]
    Assume $f: [a,b] \to \mathbb{R}$ is Riemann integrable on $[a,b]$. Show that if one value of $f$ is changed at a finite number of points in $[a,b]$, then the new function is also Riemann integrable on $[a,b]$ and the value of the integral does not change.
\end{theorem}

\clearpage

\begin{heading}
    Functions with Countable Number of Discontinuities
\end{heading}

\vspace{10pt}

\begin{theorem}
    Let $f$ be a bounded function on a closed interval $[a,b]$ that has only a countable number of discontinuities. Then $f$ is Riemann integrable on $[a,b]$.
    
\end{theorem}

\clearpage

\begin{problem}[Changing value at countable number of points]
    Assume $f: [a,b] \to \mathbb{R}$ is Riemann integrable on $[a,b]$. Show that if one value of $f$ is changed at a countable number of points in $[a,b]$, then the new function may not be Riemann integrable on $[a,b]$.
\end{problem}







\EndPage

\end{document}